% English translation (COMIMa) of the original (French) statement.
Principle of bubble sort: go through the list and swap two consecutive elements if they are not in the right order. After $n - 1$ such passes ($n$ is the length), the list is sorted.

Simplification: after the first pass, the largest value is in the last position, which is its final place; on the next pass one therefore stops at the second-to-last element, and so on.

Write a function \texttt{tri\_bulles(nombres)} that sorts the list of integers \texttt{nombres} in place with this algorithm. It returns nothing and uses no other list.

Examples: \texttt{[2, 11, 3, 7, 5]} becomes \texttt{[2, 3, 5, 7, 11]}; \texttt{[42]} and \texttt{[]} do not change.
